Актуальность темы курсового проекта обусловлена необходимостью повышения
эффективности управления заказами в различных условиях и при наличии ограниченний
складских и транспортных ресурсов, требований к срокам исполнения.
В системах логистики применение статических правил обработки заявок
может приводит к росту времени обработки заказов, увеличению количества
операций и снижение эффективности при управлении очередью заказов.

Одним из подходов к решению указанной проблемы является волновое планирование,
при котором заказы объединяются в группы по заданным критериям: приоритету,
временному окну, зоне доставки, ограничениям ресурсов и т.д.
Механизм планирования позволит увеличить устойчивость системы, 
уменьшить время задержек и повысить предсказуемость выполнения заказов.

Объектом исследования являются процессы управления заказами в системах логистики.
Предметом исследования являются методы объектно-ориентированного проектирования
системы управления и оркестрации заказов на основе волнового планирования.

Целью курсового проекта является проектирование структуры программной системы,
обеспечивающей формирование волн обработки заказов, контроль их жизненного цикла.

Для достижения поставленной цели в работе решаются следующие задачи:
\begin{itemize}
  \item выполнить анализ существующих подходов к оркестрации заказов;
  \item сформировать функциональные и нефункциональные требования к системе;
  \item провести обзор литературы по проектированию и выделить архитектурные подходы, для решения типовых архитектурных задач;
  \item разработать решения по структуре системы, доменной модели и ключевым сценариям;
  \item определить подход к тестированию и протестировать систему на наборе тестовых сценариев.
\end{itemize}

Практическая значимость работы заключается в получении архитектуры для дальнейшей 
реализации backend-сервиса оркестрации заказов,который может быть адаптирован под
требования конкретного предприятия и расширен при необходимости.